Irreducibility on the whole standing range is Lemma~\ref{lem:doubling_range} followed by
Lemma~\ref{lem:doubling_irreducible}. Everything else rests on one expansion. For $t>0$ put
$W_t(j):=(j+1)^{-t}$ at a ladder state and $W_t(s_0)=W_t(s_f):=1$. Expanding $(2j+1)^{-t}$ and
$j^{-t}$ about $(j+1)^{-t}$ and collecting the two leading terms gives, at every ladder state
$j\geq1$,
\begin{equation}\label{eq:doubling_Wdrift}
\begin{aligned}
    \mathbb E\left(W_t(X_1)\mid X_0=j\right)-W_t(j)
    &\;=\;(j+1)^{-t-1}\Bigl[\,t-c\left(1-2^{-t}\right)(j+1)^{1-s}\Bigr]+\Upsilon_t(j),\\
    |\Upsilon_t(j)|&\;\leq\;\Xi_t\Bigl((j+1)^{-t-2}+\varepsilon(j)(j+1)^{-t-1}\Bigr),
\end{aligned}
\end{equation}
with $\Xi_t<+\infty$ depending on $t$ alone: the bracket carries the drift and the remainder is
smaller by a factor $(j+1)^{-1}$ at least. Each row is then a Foster or a supermartingale
criterion applied to a suitable test function.

For $s>1$ the drift of the state itself is $c(j+1)^{1-s}-1\rightarrow-1$
(Proposition~\ref{prop:doubling_drift}), so $V_1(j):=j$ is a Foster function off a finite set
and $X$ is positive recurrent. For $s<1$ the bracket at $t=1$ tends to $-\infty$, which makes
$W_1$ a bounded non-negative supermartingale off a finite set taking somewhere off that set a
value below its infimum on it; the standard transience criterion applies.

At $s=1$ the bracket is the constant $-\phi(t)$ with $\phi(t):=c\left(1-2^{-t}\right)-t$, and
everything turns on the sign of $\phi'(0)=c\ln2-1$. For $c<1$ the return time to $s_0$ is one
step to $s_f$, one step into $\{1,\dots,d\}$ and then $\sigma$, of total expectation
$2+\bar\sigma<+\infty$ by Proposition~\ref{prop:doubling_length}: $X$ is positive recurrent.
For $c<1/\ln2$ the function $\log(j+1)$ has drift $(c\ln2-1)/(j+1)+O(j^{-2})\leq0$ off a finite
set, and Foster's recurrence criterion applies. For $c>1/\ln2$ some $t>0$ has $\phi(t)>0$, the
bracket is a negative constant, and the criterion of the case $s<1$ gives transience.

The remaining clause, $\mathbb E(\sigma\mid X_0=j)=+\infty$ at $1\leq c<2$, is an
optional-stopping estimate. The drift being $c-1\geq0$, $(X_{n\wedge\sigma})_n$ is a
submartingale; at the exit time $\varsigma_r$ of $[1,r)$ the stopped chain is bounded by
$2r-2$ and $\mathbb E(X_{\varsigma_r}\mid X_0=1)\geq1$, so the chain exits above with
probability at least $1/(2r)$, and the excursion maximum exceeds $r$ at least that often.
Summing over $r$ makes that maximum non-integrable, and
$\max_{0\leq n\leq\sigma}X_n\leq\sigma$ (Lemma~\ref{lem:doubling_excursion}) transfers this to
$\sigma$; the same holds from any $j\geq2$, which reaches $1$ before $0$ with positive
probability. An expected return time that is infinite at one state excludes positive
recurrence for an irreducible chain, and together with the recurrence just proved gives null
recurrence for $1\leq c<1/\ln2$.
